\chapter{Contexto y antecedentes}
\label{cap:contexto}

Este capítulo sitúa la pregunta de investigación en tres planos. El primero es operativo:
por qué importa caracterizar el terreno que atraviesa un rover. El segundo es histórico:
cómo ha evolucionado la medición de rocas en superficies planetarias, desde el recuento
manual hasta la segmentación automática. El tercero es analítico: qué enfoques teóricos y
metodológicos han empleado los trabajos previos, qué respuestas han obtenido y qué queda sin
resolver. El capítulo cierra con la brecha que motiva el trabajo y con las hipótesis y
variables que lo operacionalizan.

% =============================================================================
\section{Caracterización del contexto}
\label{sec:contexto}

\subsection{La operación de un rover: latencia y planificación}
\label{sec:operacion}

Un rover marciano no se conduce en tiempo real. \citet{rankin2020mobility}, que analizan las
\num{738} travesías de \textit{Curiosity} en sus siete primeros años, lo establecen como
restricción de partida: la intervención humana durante una travesía no es posible por la
latencia en la subida de comandos y la bajada de telemetría, de modo que el equipo de
operaciones depende del software de a bordo para impedir que el vehículo alcance un estado
inseguro. La comunicación entre el rover y la Tierra es, además, intermitente y diferida
\citep{tian2024rsunet}. En la práctica, cada desplazamiento se planifica desde la Tierra a
partir de las imágenes recibidas, y el rover ejecuta la secuencia con autonomía limitada.

En ese esquema, las cámaras de navegación cumplen un papel doble. Sirven al equipo en Tierra
para decidir la ruta, y sirven al propio rover para estimar su movimiento durante la
travesía: los \textit{Mars Exploration Rovers} calculaban su desplazamiento comparando pares
estéreo de sus cámaras de navegación, lo que les permitía detectar y compensar el
deslizamiento en pendientes y terreno arenoso \citep{maimone2007vo}. Las imágenes que analiza
este trabajo provienen de ese mismo tipo de instrumento.

\subsection{Las rocas como condicionante de la travesía}

La naturaleza del terreno condiciona directamente el éxito de una travesía. De los \num{738}
intentos documentados por \citet{rankin2020mobility}, \num{116} fueron impedidos o detenidos
por el software de protección, y los riesgos principales identificados para la movilidad son
el desgaste de las ruedas, su atrapamiento, el hundimiento progresivo —que puede llevar a que
el vehículo quede inmovilizado— y las interacciones con el terreno.

El caso de las ruedas de \textit{Curiosity} es ilustrativo. \citet{arvidson2017wheels}
documentan que las primeras travesías sobre afloramientos de arenisca de aristas agudas
produjeron una tasa inaceptable de perforaciones y grietas en la delgada piel de aluminio de
las ruedas: las ruedas que encabezaban un pivote de la suspensión eran empujadas contra
superficies afiladas e inmóviles por las demás. La respuesta operativa fue elaborar un mapa
geomorfológico del terreno para planificar rutas que minimizaran el daño. La lección de la
misión registra el mismo episodio \citep{jpl2017wheels}. Ese mapa es precisamente el tipo de
producto al que contribuye una caracterización sistemática del terreno: no sustituye la
decisión del equipo, pero le da información sobre dónde está la roca y cómo se organiza.

Las rocas condicionan también la fase de descenso. La selección de los sitios de aterrizaje
de las misiones a Marte ha incluido explícitamente la abundancia de rocas entre sus
restricciones de ingeniería \citep{golombek2003mer, golombek2012msl}, y el problema se ha
formalizado como el diseño de sistemas de aterrizaje frente a una distribución de obstáculos
\citep{lorenz2023analytic}.

\subsection{El conjunto de datos AI4Mars}
\label{sec:ai4mars}

AI4Mars \citep{swan2021ai4mars} es, según sus autores, el primer conjunto de datos a gran
escala para entrenar y validar modelos de clasificación de terreno en Marte. Reúne del orden
de \num{326000} etiquetas de segmentación semántica sobre unas \num{35000} imágenes de los
rovers \textit{Curiosity}, \textit{Opportunity} y \textit{Spirit}, obtenidas mediante
anotación colaborativa: cada imagen fue etiquetada por unas diez personas para asegurar el
acuerdo, y las etiquetas individuales se fusionaron exigiendo un consenso mínimo por píxel.
El conjunto incluye además del orden de \num{1500} etiquetas de validación elaboradas por
planificadores de travesías y científicos de la misión, destinadas a evaluar modelos frente a
una referencia de especialista. Los autores motivan el conjunto señalando que los vehículos
autónomos en Marte se conducían todavía con sistemas clásicos de visión artificial, y que la
segmentación semántica beneficiaría la seguridad y la productividad de las misiones. Este
trabajo emplea la versión 0.6 \citep{swan2022ai4marsdata}.

La escala de navegación distingue cuatro clases —suelo, lecho rocoso, arena y roca grande— y
un valor para el píxel sin etiqueta. La distinción entre \textit{lecho rocoso} y \textit{roca
grande} es central para este trabajo: la primera designa roca expuesta en continuidad con el
sustrato, y la segunda, bloques discretos apoyados sobre el terreno. Es una distinción
geológicamente sensata, pero, como se verá, su aplicación práctica condiciona qué puede
contarse.

% =============================================================================
\section{Antecedentes y estado actual del tema}
\label{sec:antecedentes}

La medición de rocas en superficies planetarias ha seguido una evolución en tres etapas, que
se distinguen por quién identifica las rocas y por qué producto se obtiene.

\paragraph{Primera etapa: recuento por especialistas para la seguridad de las misiones.} El
interés cuantitativo por las rocas nació de la necesidad de aterrizar. \citet{golombek1997rocks}
establecieron, a partir de las rocas medidas en las imágenes estéreo de los sitios de las
sondas \textit{Viking} y de un conjunto de análogos terrestres, que las distribuciones
tamaño--frecuencia de rocas siguen curvas exponenciales simples, y las combinaron con datos
de teledetección para predecir la frecuencia de bloques peligrosos para futuras misiones.
Ese modelo exponencial, conocido después como modelo de Golombek y Rapp, se convirtió en
referencia para la selección de los sitios de los \textit{Mars Exploration Rovers}
\citep{golombek2003mer}, del \textit{Mars Science Laboratory} \citep{golombek2012msl} y de
\textit{InSight} \citep{golombek2021insight}. En esta etapa la identificación de las rocas es
manual, se hace sobre pocas escenas y el producto es una distribución métrica.

\paragraph{Segunda etapa: detección automática para la autonomía científica.} A mediados de
la década de 2000 el objetivo se desplazó hacia la autonomía del rover: detectar rocas en las
imágenes a bordo para decidir qué observar sin esperar instrucciones desde la Tierra
\citep{castano2007oasis}. Surgieron algoritmos de detección basados en propiedades de la
imagen —bordes, sombreado, textura, forma— que \citet{thompson2007rockdet} compararon
sistemáticamente, y métodos de segmentación de rocas que combinaban rasgos a varias escalas
\citep{dunlop2007multiscale}. En paralelo, la detección de obstáculos se aplicó al descenso
\citep{huertas2006hazard}. El producto de esta etapa es la detección de instancias en
imágenes individuales, evaluada contra delimitaciones manuales.

\paragraph{Tercera etapa: segmentación semántica con aprendizaje profundo.} La disponibilidad
de conjuntos etiquetados a gran escala, y en particular de AI4Mars \citep{swan2021ai4mars},
desplazó el énfasis hacia la segmentación semántica con redes profundas: clasificar cada
píxel según el tipo de terreno. La revisión sistemática de \citet{qu2025review} recoge la
producción entre 2019 y 2025. El producto de esta etapa es un modelo, y su criterio de éxito
es el acuerdo de las predicciones con las máscaras de referencia.

\paragraph{Estado actual.} El campo combina hoy dos tradiciones que rara vez se cruzan. La de
abundancia de rocas sigue produciendo distribuciones tamaño--frecuencia rigurosas, ahora con
detección automática sobre imágenes orbitales \citep{chen2022zhurong} y con imágenes de
superficie de rovers recientes \citep{sun2024zhurong, wang2024zhurong}. La de segmentación
semántica produce arquitecturas cada vez más especializadas, evaluadas por su acuerdo con las
máscaras. La segunda trata las máscaras como la respuesta correcta; la primera no las usa. Ese
desencuentro es el que este trabajo examina.

% =============================================================================
\section{Revisión de la literatura}
\label{sec:revision}

La revisión se organiza en cinco líneas de trabajo, que corresponden a los cinco problemas
que el trabajo debe resolver: qué se ha medido sobre las rocas y con qué fundamento, cómo se
han detectado automáticamente, cómo se segmenta el terreno, cuánto puede confiarse en una
anotación colaborativa y cómo se separan instancias en contacto. Para cada línea se describen
el enfoque teórico y metodológico de los estudios y las respuestas que obtuvieron.

% -----------------------------------------------------------------------------
\subsection{Abundancia de rocas y distribuciones tamaño--frecuencia}
\label{sec:rev-abundancia}

\paragraph{Enfoque teórico.} Esta línea parte de un supuesto: que la población de rocas de una
superficie puede describirse con una función analítica de su tamaño. El objeto de estudio no
es cada roca sino la distribución acumulada —número de rocas, o fracción de área, que supera
cada diámetro—. \citet{lorenz2023analytic} revisa las funciones empleadas y encuentra un
patrón: las leyes de potencia describen bien las superficies dominadas por impactos, como la
Luna y los cuerpos pequeños, mientras que donde actúan procesos de selección eólica o fluvial,
como en Marte o Titán, la población se empobrece en los tamaños extremos y la describen mejor
funciones exponenciales y afines, entre ellas la de Golombek y Rapp. Propone además una
relación empírica entre la fracción de área cubierta por rocas y la densidad de rocas de
\SI{50}{\centi\metre}, útil para el diseño de sistemas de aterrizaje.

\paragraph{Estudios y respuestas.} \citet{golombek1997rocks} midieron las rocas en las
imágenes estéreo del campo cercano de los sitios \textit{Viking} y encontraron que sus
distribuciones, como las de diversos terrenos rocosos terrestres, tienen forma convexa en
escala logarítmica y se ajustan bien con funciones exponenciales simples. \citet{ward2005spirit}
midieron tamaño y espaciamiento de clastos en las imágenes panorámicas de \textit{Spirit} en
el cráter Gusev y hallaron que en las llanuras los clastos son menores, mejor seleccionados y
más regularmente espaciados que cerca de los bordes de cráter, lo que atribuyeron al
transporte eólico. \citet{golombek2021insight} combinaron ortoimágenes del campo cercano del
módulo \textit{InSight}, panorámicas del campo lejano e imágenes orbitales de alta resolución:
las distribuciones de cuatro zonas siguen las curvas exponenciales del modelo, con una
abundancia de roca que varía entre el \SI{0.6}{\percent} y el 3--\SI{5}{\percent} y es, en
conjunto, algo superior al \SI{3}{\percent}. En la Luna, \citet{di2016change3} analizaron la
distribución tamaño--frecuencia en el sitio de alunizaje de \textit{Chang'E-3}, y
\citet{watkins2019boulder} caracterizaron las poblaciones de bloques alrededor de seis cráteres
jóvenes, encontrando que su densidad disminuye con la edad del cráter.

La línea ha incorporado detección automática. \citet{chen2022zhurong} desarrollaron un método
basado en una red convolucional para detectar del orden de \num{10000} rocas de al menos
\SI{1.4}{\metre} en imágenes orbitales del sitio de \textit{Zhurong}, estimaron una abundancia
de roca cercana al \SI{5}{\percent} y la confirmaron con las imágenes tomadas después por el
propio rover. \citet{sun2024zhurong} calcularon la distribución a partir de las cámaras de
navegación y de terreno de \textit{Zhurong}, comparando el interior de la depresión excavada
por el aterrizaje con la superficie circundante, y \citet{wang2024zhurong} analizaron la
distribución en el mismo sitio a partir de esas cámaras.

\paragraph{Aporte y límite para este trabajo.} La línea proporciona el marco interpretativo de
una distribución tamaño--frecuencia y una vara de medida externa: si una población de rocas
contadas tiene forma decreciente, es coherente con lo que se observa en superficies
planetarias. Pero sus medidas son \textbf{métricas} —exigen estereoscopía, ortoimágenes o
escala orbital— y la identificación de cada roca la hace un especialista o un detector
validado contra especialistas. Ninguno de estos estudios emplea máscaras de segmentación
semántica como dato de entrada.

% -----------------------------------------------------------------------------
\subsection{Detección y conteo automático de rocas en imágenes de superficie}
\label{sec:rev-deteccion}

\paragraph{Enfoque teórico.} Esta línea trata la roca como un objeto que debe aislarse del
fondo a partir de las propiedades de la imagen. El fundamento es fotométrico y geométrico:
una roca proyecta sombra, tiene un contorno y una textura que difieren del suelo. El problema
de fondo, que \citet{dunlop2007multiscale} formulan con claridad, es que las rocas se prestan
mal a las técnicas de segmentación visual: presentan morfologías muy diversas y carecen de una
propiedad uniforme que las distinga del suelo.

\paragraph{Estudios y respuestas.} \citet{thompson2007rockdet} revisaron siete algoritmos de
detección de rocas de la literatura de autonomía científica y los evaluaron sobre imágenes de
los \textit{Mars Exploration Rovers}, de entornos análogos terrestres y de un simulador de
terreno marciano, mostrando que el desempeño varía con las condiciones de adquisición.
\citet{dunlop2007multiscale} propusieron combinar rasgos locales —textura—, rasgos de objeto
—sombreado y forma bidimensional— y rasgos de escena —la dirección de iluminación—: una
sobresegmentación en superpíxeles se fusiona en regiones, y un modelo aprendido de la
apariencia de las rocas puntúa cada agrupación candidata. El sistema OASIS
\citep{castano2007oasis} integró la detección de rocas en un ciclo autónomo cerrado: el rover
adquiere imágenes, encuentra rocas, analiza sus propiedades e identifica las que merecen una
observación adicional. Para el descenso, \citet{huertas2006hazard} presentaron un sistema
pasivo de detección de peligros que incluye rocas a partir de cierta altura. Más
recientemente, la detección se ha abordado con redes convolucionales en entornos análogos a
Marte \citep{furlan2019rockcnn}.

\paragraph{Aporte y límite para este trabajo.} Esta línea aporta dos ideas que el trabajo
retoma. La primera, que la separación de rocas en contacto requiere información que no está
en la silueta —el sombreado y la dirección de iluminación de \citet{dunlop2007multiscale}—;
el Capítulo~\ref{cap:resultados} vuelve sobre ella al explorar el corte por sombras. La
segunda, que el resultado depende de las condiciones de imagen, lo que obliga a validar contra
una referencia humana. Su límite es la escala: estos sistemas se evalúan en muestras reducidas
de imágenes y no producen indicadores agregados para un conjunto completo.

% -----------------------------------------------------------------------------
\subsection{Segmentación semántica del terreno marciano}
\label{sec:rev-segmentacion}

\paragraph{Enfoque teórico.} La segmentación semántica formula el problema como clasificación
por píxel: asignar a cada punto de la imagen una clase de terreno. El aprendizaje es
supervisado, y la máscara anotada hace de verdad de referencia. El criterio de evaluación
estándar es el índice de intersección sobre la unión, que mide el solapamiento entre la
región predicha y la anotada.

\paragraph{Estudios y respuestas.} El trabajo pionero en clasificación de terreno para la
planificación de travesías es SPOC \citep{rothrock2016spoc}, que aplicó aprendizaje profundo
a la clasificación del terreno en las misiones a Marte. A partir de AI4Mars
\citep{swan2021ai4mars} la producción se ha orientado a dos objetivos. El primero es la
eficiencia para cómputo a bordo: SegMarsViT \citep{dai2022segmarsvit} propone un segmentador
ligero basado en transformadores móviles, adecuado para las restricciones de procesamiento y
consumo de una nave, y LBNet \citep{wei2024lbnet} una red bilateral ligera que separa una rama
de detalle espacial de otra de información semántica. El segundo es la precisión en la
delimitación de rocas, cuya variabilidad de forma, tamaño, textura y color dificulta la tarea:
RockFormer \citep{liu2023rockformer} es el primer marco de transformador en forma de U para
segmentar rocas marcianas; MarsFormer \citep{xiong2023marsformer} combina un codificador de
transformadores jerárquicos con módulos de realce de rasgos; RSU-Net \citep{tian2024rsunet}
introduce atención en una arquitectura U-Net; otras propuestas incorporan atención híbrida
para terreno no estructurado \citep{liu2023hybrid} o información de profundidad
\citep{ma2025depthformer}. Se han propuesto también conjuntos complementarios, como GMSRI
\citep{wang2021gmsri}, que mezcla imágenes reales y sintéticas generadas por redes
adversarias y las organiza por textura y estructura espacial, y alternativas basadas en
textura para separar arena y roca \citep{alkawi2025texture}. La revisión de
\citet{qu2025review} sistematiza esta producción.

\paragraph{Aporte y límite para este trabajo.} Esta línea aporta la técnica de la
comparación con aprendizaje automático del Capítulo~\ref{cap:resultados}, y una observación
de fondo: estos trabajos segmentan \textit{píxeles de roca}, no rocas individuales. La
segmentación semántica no distingue instancias, de modo que dos rocas en contacto forman una
sola región. Y su criterio de éxito —el acuerdo con la máscara— presupone que la máscara es
correcta, un supuesto que la línea siguiente pone en cuestión.

% -----------------------------------------------------------------------------
\subsection{Calidad de la anotación colaborativa}
\label{sec:rev-anotacion}

\paragraph{Enfoque teórico.} La anotación colaborativa descansa sobre una premisa de
agregación: los errores individuales de muchos anotadores no especialistas se compensan al
combinar sus juicios, de modo que el consenso se aproxima al de un especialista. La literatura
estudia cuándo se cumple esa premisa, y qué ocurre cuando las etiquetas resultantes se usan
como referencia.

\paragraph{Estudios y respuestas.} La evidencia es en general favorable para tareas de
clasificación. \citet{snow2008cheap} encontraron un acuerdo alto entre las anotaciones de
contribuyentes no especialistas y las de especialistas en cinco tareas de lenguaje, y que
entrenar con las primeras puede ser tan eficaz como hacerlo con las segundas.
\citet{sheng2008label} mostraron que repetir el etiquetado mejora la calidad de las etiquetas
y de los modelos, aunque no siempre. En astronomía, Galaxy Zoo \citep{lintott2008galaxyzoo}
obtuvo más de cuarenta millones de clasificaciones morfológicas de casi un millón de galaxias
por parte del público, consistentes con las de astrónomos profesionales. En ecología, el
proyecto Snapshot Serengeti \citep{swanson2016citizen} agregó por mayoría las clasificaciones
de una media de veintisiete voluntarios por imagen y obtuvo un \SI{98}{\percent} de acuerdo
con un conjunto verificado por especialistas, aunque con más errores en las especies raras, y
mostró que medidas de incertidumbre del consenso permiten señalar qué imágenes requieren
revisión experta.

Dos estudios matizan ese panorama, y son los más pertinentes aquí. El primero es planetario:
\citet{robbins2014craters} compararon la identificación de cráteres lunares por ocho
investigadores profesionales y por voluntarios, y encontraron que incluso el acuerdo entre
especialistas depende del diámetro de los cráteres, de su densidad y del tipo de terreno, con
diferencias de hasta un \SI{45}{\percent} y artefactos cerca del tamaño mínimo estudiado; de
ahí su recomendación de cautela al interpretar pequeñas variaciones de las distribuciones y
de preferir el consenso de muchos anotadores. El segundo concierne a la evaluación de modelos:
\citet{northcutt2021labelerrors} estimaron que los conjuntos de prueba de diez de los conjuntos
de datos más usados en aprendizaje automático contienen en promedio al menos un
\SI{3.3}{\percent} de etiquetas erróneas, suficiente para alterar la comparación entre
modelos, y recomendaron evaluarlos sobre conjuntos de prueba correctamente etiquetados.

\paragraph{Aporte y límite para este trabajo.} Esta línea justifica examinar la anotación en
lugar de darla por buena, y ofrece el procedimiento: contrastar contra una referencia de
especialista. También advierte de una dificultad que el trabajo encontró en su propio diseño:
la comparación entre fuentes de anotación solo es interpretable si ambas etiquetan el mismo
material. Su límite para el caso presente es que los estudios favorables se refieren sobre
todo a clasificación de imágenes completas; la anotación por píxel con una regla de consenso,
como la de AI4Mars, añade un efecto —qué píxeles quedan sin etiqueta— que esos estudios no
abordan.

% -----------------------------------------------------------------------------
\subsection{Separación de instancias mediante morfología matemática}
\label{sec:rev-morfologia}

\paragraph{Enfoque teórico.} La morfología matemática trata la imagen como un objeto
geométrico y la transforma mediante operaciones definidas por un elemento estructurante
\citep{soille2003morph}. Para separar objetos en contacto, su herramienta central es la
división de aguas, que interpreta una imagen como un relieve topográfico e inunda sus cuencas
desde los mínimos; las fronteras entre cuencas definen la segmentación. \citet{vincent1991watersheds}
dieron la formulación eficiente por simulación de inmersión, y \citet{meyer1990morph}
establecieron el esquema de segmentación controlada por marcadores, en el que la inundación
parte de semillas elegidas de antemano para evitar la sobresegmentación.

\paragraph{Estudios y respuestas.} Aplicada a una imagen binaria, la división de aguas se
ejecuta sobre la transformada de distancia, cuyos máximos marcan los centros de los objetos y
cuyos mínimos locales entre máximos marcan los estrangulamientos. La selección de semillas es
el punto crítico, y la reconstrucción morfológica en escala de grises de
\citet{vincent1993reconstruction}, basada en la noción de máximos regionales, proporciona el
algoritmo eficiente sobre el que se calculan los máximos con un umbral de prominencia. El
esquema es de uso corriente donde hay que contar objetos que se tocan. En minería se aplica a
la distribución de tamaños de fragmentos tras una voladura, hoy también con redes neuronales
\citep{guo2023blasted, sharifi2024fragment}, y en cartografía de riesgo por caída de rocas se
emplea sobre inventarios obtenidos por teledetección \citep{grady2024rockfall}. La
implementación utilizada es la de \texttt{scikit-image} \citep{vanderwalt2014skimage}.

\paragraph{Aporte y límite para este trabajo.} Esta línea aporta el método del conteo y, con
él, su límite conocido: la división de aguas solo separa donde la geometría del objeto
presenta un estrangulamiento. Si la máscara no codifica la frontera entre dos objetos, ningún
procedimiento que opere únicamente sobre ella puede recuperarla.

% -----------------------------------------------------------------------------
\subsection{Síntesis comparativa}
\label{sec:rev-sintesis}

La Tabla~\ref{tab:estado-arte} resume los estudios más próximos a la pregunta de
investigación, con atención a dos rasgos que los distinguen de este trabajo: si cuentan
instancias de roca y si emplean una máscara de segmentación como dato de entrada.

\begin{landscape}
\begin{spacing}{1}
\footnotesize
\setlength{\tabcolsep}{3.5pt}
\begin{longtable}{@{}p{2.6cm}p{3.0cm}p{3.2cm}p{1.8cm}p{3.2cm}p{3.6cm}p{1.5cm}p{1.5cm}@{}}
  \caption{Síntesis comparativa de estudios previos.}
  \label{tab:estado-arte}\\
  \toprule
  \textbf{Trabajo} & \textbf{Datos} & \textbf{Objetivo} & \textbf{Unidad} & \textbf{Técnica} &
  \textbf{Resultado} & \textbf{¿Cuenta instancias?} & \textbf{¿Usa máscara?} \\
  \midrule
  \endfirsthead
  \toprule
  \textbf{Trabajo} & \textbf{Datos} & \textbf{Objetivo} & \textbf{Unidad} & \textbf{Técnica} &
  \textbf{Resultado} & \textbf{¿Cuenta instancias?} & \textbf{¿Usa máscara?} \\
  \midrule
  \endhead
  \bottomrule
  \endfoot
  \citet{golombek1997rocks} & Imágenes estéreo \textit{Viking} y análogos terrestres &
    Modelar la población de rocas & Sitio & Medición manual, ajuste de funciones &
    Distribuciones exponenciales simples & Sí, manual & No \\
  \citet{ward2005spirit} & Panorámicas de \textit{Spirit} & Tamaño y espaciamiento de clastos &
    Zona & Medición sobre imagen & Clastos menores y más espaciados en llanuras & Sí, manual & No \\
  \citet{golombek2021insight} & Ortoimágenes, panorámicas y orbitales de \textit{InSight} &
    Abundancia de rocas & Zona & Medición manual, modelo exponencial &
    Abundancia algo superior al \SI{3}{\percent} & Sí, manual & No \\
  \citet{chen2022zhurong} & Imágenes orbitales del sitio de \textit{Zhurong} &
    Abundancia y erosión & Región & Red convolucional de detección &
    \num{10000} rocas; abundancia cercana al \SI{5}{\percent}, confirmada en superficie &
    Sí, automática & No \\
  \citet{sun2024zhurong} & Cámaras de navegación y terreno de \textit{Zhurong} &
    Distribución tamaño--frecuencia & Zona & Medición sobre imagen de superficie &
    Menor abundancia en la depresión de aterrizaje & Sí & No \\
  \citet{thompson2007rockdet} & Imágenes MER, análogas y sintéticas & Comparar detectores &
    Imagen & Siete algoritmos de detección & Desempeño dependiente de la adquisición &
    Sí, automática & No \\
  \citet{dunlop2007multiscale} & Imágenes MER & Detectar y segmentar rocas & Roca &
    Superpíxeles, rasgos multiescala, modelo aprendido & Uso de sombreado e iluminación &
    Sí, automática & No \\
  \citet{castano2007oasis} & Rover FIDO en análogo & Ciencia autónoma & Roca &
    Detección y análisis a bordo & Ciclo autónomo cerrado demostrado & Sí, automática & No \\
  \citet{swan2021ai4mars} & \num{35000} imágenes de tres rovers & Conjunto de entrenamiento &
    Píxel & Anotación colaborativa con consenso & Primer conjunto a gran escala & No &
    Produce las máscaras \\
  \citet{liu2023rockformer} & Imágenes marcianas etiquetadas & Segmentar roca & Píxel &
    Transformador en U & Mejora del acuerdo con la máscara & No & Como verdad \\
  \citet{xiong2023marsformer} & Imágenes marcianas etiquetadas & Segmentar roca & Píxel &
    Transformador jerárquico & Mejora del acuerdo con la máscara & No & Como verdad \\
  \citet{dai2022segmarsvit} & Imágenes marcianas etiquetadas & Segmentación a bordo & Píxel &
    Transformador móvil ligero & Eficiencia para cómputo a bordo & No & Como verdad \\
  \citet{robbins2014craters} & Imágenes lunares & Variabilidad de anotadores & Cráter &
    Comparación de expertos y voluntarios & Acuerdo dependiente del tamaño y del terreno &
    Sí, manual & No \\
  \citet{northcutt2021labelerrors} & Diez conjuntos de referencia & Errores en conjuntos de
    prueba & Etiqueta & Aprendizaje confiable y revisión humana &
    \SI{3.3}{\percent} de errores medios; altera comparaciones & No & Audita etiquetas \\
  \textbf{Este trabajo} & \textbf{\cNEscenas{} máscaras MSL NavCam} &
    \textbf{Indicadores derivados de la máscara y sus límites} & \textbf{Imagen} &
    \textbf{Morfología matemática y división de aguas} & \textbf{Ver Capítulo~\ref{cap:resultados}} &
    \textbf{Sí, aproximado} & \textbf{Como dato} \\
\end{longtable}
\end{spacing}
\end{landscape}

La tabla deja ver la estructura del campo. Los estudios que \textbf{cuentan instancias} no usan
máscaras de segmentación: cuentan sobre la imagen, a mano o con un detector. Los que usan
máscaras las tratan \textbf{como verdad} para entrenar y evaluar modelos que clasifican
píxeles, sin contar instancias. Y los que \textbf{auditan etiquetas} lo hacen sobre
clasificación de imágenes completas o de objetos discretos, no sobre máscaras de terreno.

% =============================================================================
\section{Brecha identificada}
\label{sec:brecha}

De la revisión se desprende un espacio poco explorado. En la revisión realizada \textbf{no se
identificaron trabajos que empleen las máscaras de segmentación de AI4Mars como dato de entrada
para derivar indicadores cuantitativos del terreno}, ni que examinen los límites de lo que esas
máscaras permiten inferir. Las máscaras se han usado como respuesta para entrenar modelos, no
como medición.

Ese uso tiene interés por sí mismo y por una razón añadida. Por sí mismo, porque las máscaras
cubren más de dieciséis mil escenas de la cámara de navegación de \textit{Curiosity}, una
escala inalcanzable para la medición manual de la primera línea de la revisión. Y por la razón
añadida que sugiere la cuarta línea: si las máscaras son la referencia contra la que se evalúa
toda una familia de modelos, examinar qué contienen —qué puede medirse con ellas, qué no, y
cuánto dependen de quién y cómo las produjo— es información pertinente para esa familia de
trabajos.

El trabajo se sitúa en ese espacio. No propone una arquitectura nueva ni compite con los
modelos de segmentación: toma la máscara como dato y pregunta qué indicadores pueden
derivarse de ella de forma reproducible, y con qué límites de validez.

% =============================================================================
\section{Hipótesis y variables de estudio}
\label{sec:hipotesis}

El anteproyecto no formuló hipótesis estadísticas: planteaba la pregunta de forma descriptiva
y preveía, como validación opcional, comprobar si el conteo tiende a subestimar o a
sobreestimar frente a un observador. Para hacer contrastable el diseño se explicitan aquí las
hipótesis de trabajo que ese diseño implica. Se indica en cada caso si estaba prevista en el
anteproyecto o se incorporó durante la ejecución, para que el lector distinga lo planificado
de lo que surgió al validar.

\begin{description}[leftmargin=1.2cm, style=nextline]

\item[H1.] \textbf{La cobertura de roca visible puede derivarse de forma reproducible de las
  máscaras, y su valor no está determinado por la fracción de la escena que fue etiquetada.}
  Prevista: la cobertura es el objetivo E1; el control del denominador es la condición para
  que el indicador sea interpretable. Se contrasta con la relación entre cobertura y fracción
  etiquetada, mediante medidas de dependencia lineal y no lineal.

\item[H2.] \textbf{El conteo por división de aguas sobre la clase roca grande aproxima el número
  de rocas que un observador distingue dentro de la región anotada, sin sesgo sistemático en
  una dirección.} Prevista: es el objetivo E2, y el anteproyecto contemplaba comprobar la
  dirección del error. Se contrasta con el acuerdo por bandas frente al conteo humano, medido
  con kappa simple y ponderado, y con la proporción de escenas por encima y por debajo.

\item[H3.] \textbf{Los indicadores derivados no dependen de la fuente de la anotación.}
  Incorporada durante la ejecución, al constatar que el conjunto incluye máscaras de experto.
  El anteproyecto daba las etiquetas colaborativas por fiables. Se contrasta comparando ambas
  fuentes y separando, mediante un instrumento común, la parte de la diferencia atribuible a
  las imágenes de la atribuible a la anotación.

\item[H4.] \textbf{Cobertura y conteo aportan información distinta sobre el terreno.} Implícita
  en el diseño de dos indicadores: si uno fuera función del otro, bastaría con uno. Se
  contrasta con medidas de dependencia y con la fracción de la varianza de cada indicador que
  explica el otro.

\end{description}

La Tabla~\ref{tab:variables} define las variables de estudio. Todas se derivan de la máscara,
excepto la banda del observador, que se recoge en la validación, y la cobertura del modelo, que
procede del segmentador entrenado.

\begin{table}[H]
  \centering
  \small
  \caption{Variables de estudio.}
  \label{tab:variables}
  \begin{tabular}{p{3.6cm}p{6.4cm}p{2.6cm}c}
    \toprule
    \textbf{Variable} & \textbf{Definición operativa} & \textbf{Tipo} & \textbf{Hip.} \\
    \midrule
    Cobertura sobre lo etiquetado & Píxeles de lecho rocoso o roca grande sobre píxeles
      etiquetados, en porcentaje & Continua, 0--100 & H1, H3, H4 \\
    Cobertura sobre la imagen & Los mismos píxeles de roca sobre el total de la imagen &
      Continua, 0--100 & H1 \\
    Fracción etiquetada & Píxeles etiquetados sobre el total de la imagen & Continua, 0--1 & H1 \\
    Número de rocas & Regiones de roca grande tras la división de aguas que superan los filtros
      de área y forma & Discreta, $\geq 0$ & H2, H4 \\
    Componentes conectadas & Regiones de roca grande antes de la división de aguas &
      Discreta, $\geq 0$ & H2 \\
    Composición & Proporción de cada clase sobre lo etiquetado & Continua, 0--100 & H3 \\
    Banda del observador & Número de rocas distinguidas por una persona dentro de la región
      anotada, en seis bandas & Ordinal & H2 \\
    Fuente de la anotación & Colaborativa o de especialista & Nominal & H3 \\
    Cobertura del modelo & Fracción de roca que predice el segmentador en la región anotable
      de la escena & Continua, 0--100 & H3 \\
    Reloj de nave & Contador temporal de a bordo, extraído del identificador & Ordinal & Secuencia \\
    \bottomrule
  \end{tabular}
\end{table}
